\documentclass[10pt]{article}
\usepackage[margin=20mm]{geometry}
\usepackage{amsmath,amssymb,fontspec,unicode-math,xurl,hyperref}
\IfFontExistsTF{DejaVu Serif}{\setmainfont{DejaVu Serif}}{\setmainfont{Latin Modern Roman}}
\setmathfont{latinmodern-math.otf}
\setlength{\parindent}{0pt}\setlength{\parskip}{6pt}\sloppy
\title{Technical Verification: LOGO – Long cOntext aliGnment via efficient preference Optimization}
\author{AI-assisted technical review}
\date{}
\begin{document}
\maketitle
Abstract: An exact bounded-probability counterexample invalidates the generalization bound used to justify LOGO theoretically. Independent derivative and finite-tail calculations establish the defect while leaving the reported empirical improvements unrefuted. Static inspection also identifies a two-negative slicing error in the archived training release and a written SFT sign discrepancy. Historical experimental consequences of the code issue are unresolved. Scientific adjudication remains pending.\par
\section*{Paper identity and reviewed version}
This report concerns LOGO – Long cOntext aliGnment via efficient preference Optimization, by Zecheng Tang et al., ICML 2025, PMLR volume 267; identifier `I25-43672`, OpenReview `vVEBtDDSA6`. The reviewed object is the frozen published proceedings PDF, retrieved on 25 October 2026. Page references below are one-based PDF pages, which coincide with the internal page numbers. They are not the continuous PMLR volume pagination.\par
\section*{Scope and scientific claims}
LOGO combines synthetic positional indices with preference optimization over one preferred and multiple rejected responses. This review covers the main method and experiments, the generalization argument in Appendix C (pages 14–16), the gradient discussion in D, the positional argument in E, and the experimental extensions and examples in F–K. The publicly archived repository was read statically at commit `a960d8496ff56108804cf36c8bcfc7b039236c92`; it is a retrieved release, not established as the exact code used for the published experiments. Source pages 5, 8, 15 and 16 were visually inspected. Neither model training nor benchmark evaluation was reproduced.\par
Section 3.2.2, page 5, expressly claims that Appendix C proves a small generalization error and reliable performance on unseen data. Its per-example preference loss is \[\ell(z)=-\log\sigma(z)=\log(1+e^{-z}),\quad
z=\frac\beta{|y_w|}\log\pi_\theta(y_w\mid x)-\frac\beta{M|y_l|}\sum_{j=1}^M\log\pi_\theta(y_l^{(j)}\mid x)-\gamma.\] The following calculation grants the appendix’s strong bounded-log-probability assumption; its failure does not rely on questioning that assumption.\par
\section*{A false concentration guarantee}
Problem 1 (Appendix C’s Lipschitz and bounded-difference calculation fails).  The asserted Lipschitz constants and resulting generalization probability bound are invalid, including on a finite, bounded-probability example with a fixed model and independent samples.\par
Equation reference: eq:claimed. Step 4 on page 15 calls \(\ell\) \(1/4\)-Lipschitz and displays \(|\ell'(z)|=1/(1+e^z)\leq1/2\) for every real \(z\). In fact this derivative magnitude exceeds \(1/2\) whenever \(z<0\) and approaches one as \(z\to-\infty\). The globally valid Lipschitz constant is one. The constant \(1/4\) is associated with the sigmoid derivative, not its negative logarithm. This is not a harmless label: page 16 uses \(c_i=|z_{\rm new,i}-z_{\rm old,i}|/2\) in \[\Pr\{|L-\widehat L|\geq\epsilon\}\leq
2\exp\left(-\frac{2N\epsilon^2}{\sum_{i=1}^N c_i^2}\right).\]\par
Here is a fully specified counterexample to this displayed bound. Set \(M=2\), \(\beta=\gamma=1\), and all response lengths to one. There are two possible inputs \(x_A,x_B\). The four response probabilities, in the order preferred, rejected 1, rejected 2, remaining, are \[\begin{aligned}
\pi(\cdot\mid x_A)&=(e^{-3},e^{-2},e^{-2},1-e^{-3}-2e^{-2}),\\
\pi(\cdot\mid x_B)&=(e^{-2},e^{-2},e^{-2},1-3e^{-2}).\end{aligned}\] These are normalized strictly positive distributions and every probability is at least \(e^{-3}\). Thus \(C=3\) satisfies the assumption on this response space. One can take the parameter space to consist of this single fixed model, so there is no adaptive-training issue in the counterexample. The preferred/rejected labels need not coincide with the model’s probability ranking. The resulting scores are \(z_A=-2\), \(z_B=-1\).\par
Let the data distribution give probabilities \(0.9\) and \(0.1\) to these two tuples, respectively, and take \(N=1\). The losses are \(2.1269280110\) and \(1.3132616875\), with range \(0.8136663235\). The maximum score replacement difference is one, so the paper’s choice is \(c_1=1/2\). For \(\epsilon=0.7\), observing \(x_B\) gives \(L-\widehat L=0.9(0.8136663235)=0.7322996912\); observing \(x_A\) gives absolute deviation \(0.0813666324\). Hence the exact two-sided tail probability is \(0.1\), whereas [eq:claimed] (\#eq:claimed) gives \[2e^{-2(0.7)^2/(0.5)^2}=0.0396821895<0.1.\] An independently written script checks normalization, the probability floor, both outcomes and the numerical inequality. No simulation or author code is involved. The claim is a general bound for all \(N\); one admissible \(N\) suffices to disprove it. This does not imply that the experiments at their much larger training sample size have poor generalization.\par
There is a separate normalization problem. For a fixed model with loss range \(R\), changing one sample changes the empirical mean by at most \(R/N\). Bounded differences then gives \(2\exp(-2N\epsilon^2/R^2)\). If \(c_i\) instead denotes the change in the unaveraged example loss, the numerator in the paper’s form requires \(N^2\), not \(N\). Indeed, with its own constant \(c_i=1/2\), the printed exponent is independent of \(N\), contrary to the stated tightening with increasing sample size.\par
Replacing the constants and normalization proves a fixed-model concentration statement only. A trained parameter depending on these same samples additionally needs an appropriate uniform-complexity or stability argument. None is supplied in Appendix C. Nor does the proof assign constants and a confidence level that establish that the reported 6,000 source examples or 0.3 billion tokens suffice. The \(N=1\) example is deliberately independent of that additional problem.\par
The supported assessment is Major: the proof of a prominently claimed theoretical guarantee is false. The empirical LongBench, retrieval and language-modeling observations are separate evidence and are not refuted. A corrected fixed-model bound is verified above; a guarantee for the actually trained neural model and the claimed sample sufficiency remains unverified.\par
\section*{Objective and released implementation discrepancies}
Problem 2 (The written SFT regularizer has the opposite sign).  Equation (3), page 5, adds \(+\lambda\mathbb E\log\pi_\theta(y_w\mid x)\) to a minimized preference loss and specifies \(\lambda=0.1\). Minimizing this term by itself lowers the preferred-response log probability. Standard SFT uses its negative.\par
The released `training/logo\_train.py`, lines 271–274, adds positive-weight cross entropy, which has the conventional negative-log-likelihood sign. Figure 5(a), page 8, also reports improved perplexity with stronger SFT regularization, supporting an intended conventional sign. These are reasons to treat this as a Minor written-objective discrepancy, not evidence that the experiments optimized the displayed wrong-sign term. Correcting the sign resolves this algebraic issue; the full causal-label alignment of the released SFT implementation is not certified here.\par
Problem 3 (The released two-negative slicing changes the preference loss).  The instantiated `LOGOTrainer` concatenates chosen, reject-1 and reject-2 batches, but lines 245–247 slice the three log-probability blocks as `[:B]`, `[B:2*B]` and `[B:]`. The last block has \(2B\) entries instead of \(B\); line 266 then averages the two rejected blocks.\par
This is in the executed class selected by the entry point, not merely an unused legacy class. The three supplied launch scripts set per-device batch size one. At \(B=1\), broadcasting makes the two rejected averages \(r_1\) and \((r_1+r_2)/2\), so the mean loss mixes a single-negative term with the intended two-negative term. For example, independent scalar-array algebra with chosen score \(-1\), rejected scores \(-2,-4\), \(\beta=1\), \(\gamma=0\) yields \(0.22009485\) instead of the intended \(0.12692801\). This is a diagnostic parameter choice, not a claimed reproduction of the launch configuration. At \(B=2\), arrays of lengths two and four cannot be added. The matching logits slice already starts at `2*B`, further supporting a log-probability slicing mistake. The local repair `all\_logps[2*B:]` is shape- and algebra-verified, but has not been run in the trainer.\par
This is Major for fidelity of the archived released implementation to the core objective; the impact on the paper’s historical experiments is Unverifiable without their exact code and logs. The archive commit is dated 10 October 2025, so this review does not silently equate it with the earlier experiment version. No author program was executed. A corrected slice alone does not validate the full training pipeline.\par
\section*{Empirical checks and boundaries}
The data-token arithmetic \(12{,}000\cdot512\cdot16\cdot3=294{,}912{,}000\) is consistent with the rounded 0.3 billion claim. Figure 5 visually supports a compute/perplexity/performance trade-off, rather than one setting dominating all three axes. Appendix G.3 explicitly says more tokens are required for longer contexts, and Table 2 reports 0.3/1.2/4.8 billion tokens at 80K/128K/256K. Table 3’s doubled training time and reduced throughput with context parallelism provide useful empirical counterweight to overly broad efficiency readings. The qualitative cases are illustrative and do not estimate error frequencies. Appendix D’s gradient discussion and E’s positional-coverage argument were read, but not established as global convergence or full long-context-distribution equivalence theorems. No new finding is asserted from those broader limitations alone.\par
\section*{Conclusion}
The demonstrated Major issue is a false theoretical generalization bound. A separate core-objective bug is established in the archived release; its historical experimental impact is unresolved. The SFT sign is a bounded written discrepancy with contrary implementation evidence. The finite counterexample and local algebraic repairs are verified; model retraining, a learned-model generalization theorem, and author adoption remain unverified. Focused independent AI acceptance has checked all three Problems; scientific adjudication is pending and no human review has occurred.\par
1 Zecheng Tang et al., LOGO – Long cOntext aliGnment via efficient preference Optimization, PMLR 267 (2025), published proceedings version. https://proceedings.mlr.press/v267/tang25j.html.\par
\section*{Scope and limitations}
Source code supplied by the authors was not executed. The calculation is an independent finite or analytic diagnostic, not a reproduction of the original training experiments. A focused check revisited the stated Problems against the frozen original source, assumptions, decisive arguments, counterevidence and impact boundaries. No human review has occurred; scientific adjudication is pending and author adoption is unknown.\par
\end{document}
